% =====================================================================
%  multimachine.tex  --  PHASE IV, Task 8
%
%  The eight-machine characterisation: where the state definition holds,
%  where it fails, how it says so, and which machines have a decision
%  problem at all.
%
%  Intended placement: Section 8 (\label{sec:multimachine}), after the
%  decision layer.
%
%  Drop-in usage: % =====================================================================
%  multimachine.tex  --  PHASE IV, Task 8
%
%  The eight-machine characterisation: where the state definition holds,
%  where it fails, how it says so, and which machines have a decision
%  problem at all.
%
%  Intended placement: Section 8 (\label{sec:multimachine}), after the
%  decision layer.
%
%  Drop-in usage: % =====================================================================
%  multimachine.tex  --  PHASE IV, Task 8
%
%  The eight-machine characterisation: where the state definition holds,
%  where it fails, how it says so, and which machines have a decision
%  problem at all.
%
%  Intended placement: Section 8 (\label{sec:multimachine}), after the
%  decision layer.
%
%  Drop-in usage: % =====================================================================
%  multimachine.tex  --  PHASE IV, Task 8
%
%  The eight-machine characterisation: where the state definition holds,
%  where it fails, how it says so, and which machines have a decision
%  problem at all.
%
%  Intended placement: Section 8 (\label{sec:multimachine}), after the
%  decision layer.
%
%  Drop-in usage: \input{multimachine}
%  Figures are referenced by bare name; add to the preamble
%      \graphicspath{{../outputs/phase4/figures/}{../outputs/phase6/figures/}}
%  Bibliography: paper/references.bib
%  Numbers below are generated by  python run_phase4.py --only task8
%  and stored in outputs/phase4/task8/table_*.csv, with the intervals
%  from outputs/phase5/task11/state_block_bootstrap.csv.
% =====================================================================

\section{The Eight Machines}
\label{sec:multimachine}

Everything to this point is pelletizer~I. This section applies the
identical procedure --- one preprocessing, one labeller, $k=4$ fixed,
one physics battery --- to all eight IMDELD machines, and asks two
questions of the result: on which machines is the state definition
physically valid, and on which is there a shutdown decision to make at
all. The two answers are not the same, and neither follows the
prediction that motivated the experiment.

\subsection{Uniform labelling}
\label{sec:multimachine:labelling}

The Phase~I--III labels exist only for pelletizer~I and were produced by
a script whose configuration has since moved, so comparing that export
against seven fresh labellings would confound machine differences with
procedure differences. All eight records are therefore re-labelled
identically. Three departures from the original implementation were
required, each verified rather than assumed.

\emph{Blank currents.} The two milling-machine meters leave the current
channel empty on $36.1\%$ and $36.3\%$ of rows --- matching their OFF
shares almost exactly, i.e.\ the meter blanks the channel whenever the
machine is de-energised. On every such row the apparent power is exactly
$0$\,VA with voltage present, and $S=VI$ then admits only $I=0$\,A, so
the cells are filled with zero and a unit test asserts that every filled
cell met that condition. Leaving them missing makes a GMM
unfittable; dropping them would delete a third of each record precisely
where the machine is off, biasing every state-time total.

\emph{Timestamps and storage.} IMDELD spans a daylight-saving change, so
the timestamp column carries mixed UTC offsets; splitting the naive part
from the offset is exact and ten times faster than the fallback parse.
Records are stored as \texttt{float32}, whose one measurable consequence
is that about two rows per $10^{5}$ fall on the other side of the
spike threshold.

On pelletizer~I --- the one machine both labellings cover --- the
uniform labeller reproduces the pipeline's own export at $98.88\%$ row
agreement, $99.84\%$ on STANDBY-versus-rest, ARI $0.978$, and $89$\,h of
STANDBY against the export's $87$\,h. A number in this section that
differs from an earlier one is therefore a consequence of the analysis,
not of the labels.

Both of those figures are single-seed. The total this paper quotes for
pelletizer~I is the physics-conditional mean over passing seeds,
$92\pm3$\,h (Section~\ref{sec:method:accept_reject}); the single-seed
values sit inside that interval, and the per-machine rows of
Table~\ref{tab:multimachine} are reported per covered day with a
sampling interval for the same reason.

\subsection{Where the state definition holds}
\label{sec:multimachine:states}

\begin{table*}[!t]
\centering
\caption{Per-machine characterisation. STANDBY hours per covered day
carry a $95\%$ one-week moving-block bootstrap interval
(Section~\ref{sec:significance}); $n_{b}$ is the number of blocks the
record supports. $P_{\text{sb}}$ is the mean power of the inferred
STANDBY component. The no-load current ratio has an acceptance band of
$[0.25,0.50]$ from induction-motor theory. M1 is the fraction of
energised rows on which the independent Otsu power-factor classifier
agrees with the state model.}
\label{tab:multimachine}
\small
\begin{tabular}{@{}lrrrrrrrr@{}}
\toprule
Machine & cov.\ (d) & STANDBY (h/d) & $n_{b}$ & $P_{\text{sb}}$ & physics & PF sep. & $I$ ratio & M1 \\
\midrule
\multicolumn{9}{@{}l}{\textit{Motor-driven --- the state definition applies}}\\
Pelletizer I    & $63.4$ & $1.40$\,{\footnotesize$[1.15,1.63]$} & $10$ & $4.2$\,kW & $\mathbf{5/5}$ & $0.72$ & $0.36$ & $99.4\%$ \\
Pelletizer II   & $62.4$ & $1.85$\,{\footnotesize$[1.47,2.28]$} & $9$  & $4.5$\,kW & $\mathbf{5/5}$ & $0.71$ & $0.36$ & $99.7\%$ \\
Milling I       & $12.2$ & $5.52$\,{\footnotesize$[4.58,6.23]$} & $\mathbf{2}$ & $4.6$\,kW & $\mathbf{5/5}$ & $0.61$ & $0.43$ & $98.1\%$ \\
Milling II      & $12.2$ & $5.27$\,{\footnotesize$[4.44,5.89]$} & $\mathbf{2}$ & $5.4$\,kW & $\mathbf{5/5}$ & $0.52$ & $0.33$ & $97.7\%$ \\
\addlinespace
\multicolumn{9}{@{}l}{\textit{Auxiliary --- no energised-idle state exists to find}}\\
Exhaust fan I   & $63.7$ & $8.73$\,{\footnotesize$[7.31,10.03]$} & $10$ & $0.3$\,W & $4/5$ & $0.91$ & $\mathbf{0.008}$ & $96.6\%$ \\
Exhaust fan II  & $62.3$ & $1.45$\,{\footnotesize$[1.09,2.07]$}  & $9$  & $0.2$\,W & $3/5$ & $0.62$ & $\mathbf{0.002}$ & $\mathbf{55.1\%}$ \\
Contactor I     & $63.7$ & $0.90$\,{\footnotesize$[0.71,1.12]$}  & $10$ & $0.3$\,W & $3/5$ & $0.31$ & $\mathbf{0.009}$ & $90.3\%$ \\
Contactor II    & $62.5$ & $1.15$\,{\footnotesize$[0.85,1.42]$}  & $9$  & $9.3$\,W & $3/5$ & $\mathbf{-0.05}$ & $\mathbf{0.007}$ & $85.4\%$ \\
\bottomrule
\end{tabular}
\end{table*}

\textbf{The two pelletizers and the two milling machines pass all five
physics checks.} Their STANDBY components sit at $4.2$--$5.4$\,kW with a
power factor of $0.11$--$0.17$ against $0.64$--$0.65$ when loaded, and a
no-load current ratio inside the motor band $[0.25,0.50]$. That is the
textbook signature of an energised induction motor drawing magnetising
current, recovered independently on four machines from unlabelled data,
and it is the strongest evidence in this paper that the state definition
is physical rather than fitted.

\textbf{On the fans and contactors there is no energised-idle state to
find, and the battery says so.} Their ``STANDBY'' components sit at
$0.2$--$9.3$\,W drawing $0.09$--$0.14$\,A: the model is splitting the
OFF band in two because it was told to find four states. The current
ratio is $0.002$--$0.009$ against an acceptance band starting at
$0.25$ on all four, and contactor~II additionally shows a
\emph{negative} power-factor separation. \textbf{This is the physics
battery doing the job it was built for.} The framework's state
definition is specific to motor-driven loads, and on the machines where
it does not apply it reports that rather than returning four confident
states regardless.

Two mechanisms behind those failures are worth naming, because both are
limitations of the method rather than of the data.

\emph{Ranking states by active power can misname them.} Exhaust fan~I's
``OFF'' component ($0.8\%$ of the record) carries $6.6$\,A at $0$\,W and
power factor $0$ --- standing reactive current with no real power ---
while its ``STANDBY'' component ($36.4\%$) sits at $0.3$\,W and
$0.13$\,A and is the machine's true de-energised state. The canonical
names are assigned by mean active power, so a near-zero-power component
carrying reactive current ranks below the genuine off state and takes
its name. No hour total changes, since both are non-productive, but it
is why the current-ratio check fails in the direction it does. The same
machine returns the largest forbidden-transition probability in the
plant ($3.4\times10^{-3}$, against $1.9\times10^{-4}$ on the four motor
machines), which is an independent measurement arriving at the same
conclusion (Section~\ref{sec:explain:transitions}).

\emph{The independent classifier fails where the model has nothing to
lock onto.} Otsu-thresholded power factor agrees with the state model on
$85$--$99.7\%$ of energised rows on seven machines and on $\mathbf{55.1\%}$
on exhaust fan~II --- the machine whose physics checks fail hardest.
Method~1's usefulness as an external check is thereby confirmed
\emph{and} its failure mode identified.

\subsection{Two results that generalise}
\label{sec:multimachine:general}

\textbf{The minimum-dwell result holds on all eight records.} Flicker
before the constraint ranges from $26.7\%$ (milling~II) to $77.7\%$
(exhaust fan~II), median $50.6\%$; after it, $0.00\%$ on seven machines
and $0.17\%$ on milling~I, at a cost of relabelling $0.15$--$6.1\%$ of
rows. What Section~\ref{sec:results:task4} established on one machine
--- that a transition prior cannot control flicker and an explicit
dwell constraint can --- now holds across duty cycles differing by an
order of magnitude.

\textbf{BIC does not select a state count here.} Widening the scan to
$k\in[2,8]$ returns $k=8$ on seven machines and $k=7$ on pelletizer~II,
with the criterion still falling at the boundary in seven of eight
cases; the improvement from $k=4$ to $k=8$ is $3.1$--$33.1\%$ of the
criterion's own value. This is a result about the criterion, not the
data (Section~\ref{sec:rw:states}). Fixing $k=4$ everywhere for
comparability is a choice, and this is the honest statement of what it
costs: \textbf{the physics vetoes are not a belt-and-braces addition on
top of BIC --- on records of this size they are what is actually doing
the selecting.}

\subsection{Which machines are hard, and why}
\label{sec:multimachine:difficulty}

Difficulty is measured on three axes: separability of the no-load/load
boundary (Cohen's $d$ and Bhattacharyya distance on log power), record
quality (coverage, spike rate, segment count), and agreement with the
independent classifier.

The ordering is not the one predicted. The expectation was that the
low-power auxiliary machines would be hardest. What the data show is
that the fans and contactors have the \emph{largest} separability
figures in the plant --- Cohen's $d$ of $9.3$ to $75.0$ against the
pelletizers' $4.3$ to $7.2$ --- while failing the physics checks,
because the separation being measured there is OFF against running, not
no-load against loaded. \textbf{Separability and validity are close to
orthogonal on this data}, which is the argument for keeping an external
physics battery rather than an internal cluster-quality score:
silhouette, BIC and Cohen's $d$ are all maximised by the partitions that
are physically wrong.

The two milling machines are hard for an unrelated reason. Their records
cover $12.2$ of $43$ days in $1\,082$ and $580$ fragments, against
$66$--$74$ fragments for the five-month records. They support two
one-week bootstrap blocks, which is not enough to form a usable
interval, and the intervals in Table~\ref{tab:multimachine} are printed
for those rows only to make that visible.
\textbf{Their $5.5$\,h/day --- the largest per-day standby figure among
the motor-driven machines --- must not be annualised}, and
Figure~\ref{fig:standby-hours} prints the block count inside each bar so
the figure cannot be read as the largest opportunity in the plant.

\begin{figure}[t]
  \centering
  \includegraphics[width=\linewidth]{fig_p6_standby_hours}
  \caption{STANDBY hours per covered day with $95\%$ weekly-block
  intervals. Bars are marked according to whether the machine passes the
  physics battery, and the number of weekly blocks is printed inside
  each bar; hatching marks an interval that rests on too few blocks to
  be usable.}
  \label{fig:standby-hours}
\end{figure}

\subsection{The decision layer exists on three machines of eight}
\label{sec:multimachine:decision}

A decision problem requires a measurable standby power and an observable
cold restart. Five machines fail an explicit degeneracy test rather than
an inspection: four because standby power is below $1$\,W, so the
standby cost per second is numerically zero, the break-even point is
infinite, and every currency figure would be an artefact of dividing by
zero; and pelletizer~II because in five months it is \emph{never once
observed resuming production from OFF}, so the cold-versus-warm contrast
cannot be measured at all. Two machines return a \emph{negative} restart
delay --- cold starts reaching production faster than warm ones --- which
is not a plausible physical reading and is a further symptom of the same
problem. The three machines that survive the test are the subject of
Table~\ref{tab:multimachine:decision}.

\begin{table}[t]
\centering
\caption{Machines on which the shutdown decision is well posed.
Head-room is the oracle saving; the proposed policy's figure is quoted
against the plant's own behaviour.}
\label{tab:multimachine:decision}
\small
\setlength{\tabcolsep}{4pt}
\begin{tabular}{@{}lrrrr@{}}
\toprule
Machine & energised idle & head-room & forecast MAE & proposed \\
        & (h, \% of idle) & (USD/yr) & (s) & (USD/yr) \\
\midrule
Pelletizer I & $88$ ($18.5\%$) & $37.6$ & $9\,034$ & $+2.9$ \\
Milling I    & $67$ ($48.8\%$) & $24.9$ & $18\,205$ & $\mathbf{-1\,041}$ \\
Milling II   & $64$ ($47.6\%$) & $25.2$ & $18\,440$ & $\mathbf{-1\,256}$ \\
\bottomrule
\end{tabular}
\end{table}

\textbf{The conjecture that motivated this experiment is refuted.} The
open question after the single-machine study was whether some other
IMDELD machine has enough recoverable standby to fund a retrofit ---
higher standby power, a larger energised-idle fraction, or a plant that
does not already shut down manually. None does. The three machines with
a well-posed problem have head-rooms of $37.6$, $24.9$ and
$25.2$\,USD/yr, so \textbf{the facility's total attainable saving from
standby recovery, summed over every machine where the question is even
meaningful, is under $90$\,USD/yr}. The framework's contribution on this
facility is decision quality, not recovered energy, and this paper says
so in its abstract rather than at the end of its discussion.

\textbf{The milling machines are worse than that, and the negative
result is reported as it stands.} The proposed policy \emph{loses}
$1\,041$ and $1\,256$\,USD/yr there against a head-room of
$25$\,USD/yr. The cause is visible in the same table: a duration
forecaster with MAE $18\,200$\,s and Spearman $0.34$--$0.42$, against
$9\,034$\,s and $0.58$ on pelletizer~I, fitted on a record broken into
$1\,082$ fragments. It over-predicts remaining idle time and fires
shutdowns the realised episode does not support. The decision layer
inherits its forecaster's quality, and on a record this fragmented the
forecaster is not good enough to act on --- which is a deployment
criterion, not merely an observation.

%  Figures are referenced by bare name; add to the preamble
%      \graphicspath{{../outputs/phase4/figures/}{../outputs/phase6/figures/}}
%  Bibliography: paper/references.bib
%  Numbers below are generated by  python run_phase4.py --only task8
%  and stored in outputs/phase4/task8/table_*.csv, with the intervals
%  from outputs/phase5/task11/state_block_bootstrap.csv.
% =====================================================================

\section{The Eight Machines}
\label{sec:multimachine}

Everything to this point is pelletizer~I. This section applies the
identical procedure --- one preprocessing, one labeller, $k=4$ fixed,
one physics battery --- to all eight IMDELD machines, and asks two
questions of the result: on which machines is the state definition
physically valid, and on which is there a shutdown decision to make at
all. The two answers are not the same, and neither follows the
prediction that motivated the experiment.

\subsection{Uniform labelling}
\label{sec:multimachine:labelling}

The Phase~I--III labels exist only for pelletizer~I and were produced by
a script whose configuration has since moved, so comparing that export
against seven fresh labellings would confound machine differences with
procedure differences. All eight records are therefore re-labelled
identically. Three departures from the original implementation were
required, each verified rather than assumed.

\emph{Blank currents.} The two milling-machine meters leave the current
channel empty on $36.1\%$ and $36.3\%$ of rows --- matching their OFF
shares almost exactly, i.e.\ the meter blanks the channel whenever the
machine is de-energised. On every such row the apparent power is exactly
$0$\,VA with voltage present, and $S=VI$ then admits only $I=0$\,A, so
the cells are filled with zero and a unit test asserts that every filled
cell met that condition. Leaving them missing makes a GMM
unfittable; dropping them would delete a third of each record precisely
where the machine is off, biasing every state-time total.

\emph{Timestamps and storage.} IMDELD spans a daylight-saving change, so
the timestamp column carries mixed UTC offsets; splitting the naive part
from the offset is exact and ten times faster than the fallback parse.
Records are stored as \texttt{float32}, whose one measurable consequence
is that about two rows per $10^{5}$ fall on the other side of the
spike threshold.

On pelletizer~I --- the one machine both labellings cover --- the
uniform labeller reproduces the pipeline's own export at $98.88\%$ row
agreement, $99.84\%$ on STANDBY-versus-rest, ARI $0.978$, and $89$\,h of
STANDBY against the export's $87$\,h. A number in this section that
differs from an earlier one is therefore a consequence of the analysis,
not of the labels.

Both of those figures are single-seed. The total this paper quotes for
pelletizer~I is the physics-conditional mean over passing seeds,
$92\pm3$\,h (Section~\ref{sec:method:accept_reject}); the single-seed
values sit inside that interval, and the per-machine rows of
Table~\ref{tab:multimachine} are reported per covered day with a
sampling interval for the same reason.

\subsection{Where the state definition holds}
\label{sec:multimachine:states}

\begin{table*}[!t]
\centering
\caption{Per-machine characterisation. STANDBY hours per covered day
carry a $95\%$ one-week moving-block bootstrap interval
(Section~\ref{sec:significance}); $n_{b}$ is the number of blocks the
record supports. $P_{\text{sb}}$ is the mean power of the inferred
STANDBY component. The no-load current ratio has an acceptance band of
$[0.25,0.50]$ from induction-motor theory. M1 is the fraction of
energised rows on which the independent Otsu power-factor classifier
agrees with the state model.}
\label{tab:multimachine}
\small
\begin{tabular}{@{}lrrrrrrrr@{}}
\toprule
Machine & cov.\ (d) & STANDBY (h/d) & $n_{b}$ & $P_{\text{sb}}$ & physics & PF sep. & $I$ ratio & M1 \\
\midrule
\multicolumn{9}{@{}l}{\textit{Motor-driven --- the state definition applies}}\\
Pelletizer I    & $63.4$ & $1.40$\,{\footnotesize$[1.15,1.63]$} & $10$ & $4.2$\,kW & $\mathbf{5/5}$ & $0.72$ & $0.36$ & $99.4\%$ \\
Pelletizer II   & $62.4$ & $1.85$\,{\footnotesize$[1.47,2.28]$} & $9$  & $4.5$\,kW & $\mathbf{5/5}$ & $0.71$ & $0.36$ & $99.7\%$ \\
Milling I       & $12.2$ & $5.52$\,{\footnotesize$[4.58,6.23]$} & $\mathbf{2}$ & $4.6$\,kW & $\mathbf{5/5}$ & $0.61$ & $0.43$ & $98.1\%$ \\
Milling II      & $12.2$ & $5.27$\,{\footnotesize$[4.44,5.89]$} & $\mathbf{2}$ & $5.4$\,kW & $\mathbf{5/5}$ & $0.52$ & $0.33$ & $97.7\%$ \\
\addlinespace
\multicolumn{9}{@{}l}{\textit{Auxiliary --- no energised-idle state exists to find}}\\
Exhaust fan I   & $63.7$ & $8.73$\,{\footnotesize$[7.31,10.03]$} & $10$ & $0.3$\,W & $4/5$ & $0.91$ & $\mathbf{0.008}$ & $96.6\%$ \\
Exhaust fan II  & $62.3$ & $1.45$\,{\footnotesize$[1.09,2.07]$}  & $9$  & $0.2$\,W & $3/5$ & $0.62$ & $\mathbf{0.002}$ & $\mathbf{55.1\%}$ \\
Contactor I     & $63.7$ & $0.90$\,{\footnotesize$[0.71,1.12]$}  & $10$ & $0.3$\,W & $3/5$ & $0.31$ & $\mathbf{0.009}$ & $90.3\%$ \\
Contactor II    & $62.5$ & $1.15$\,{\footnotesize$[0.85,1.42]$}  & $9$  & $9.3$\,W & $3/5$ & $\mathbf{-0.05}$ & $\mathbf{0.007}$ & $85.4\%$ \\
\bottomrule
\end{tabular}
\end{table*}

\textbf{The two pelletizers and the two milling machines pass all five
physics checks.} Their STANDBY components sit at $4.2$--$5.4$\,kW with a
power factor of $0.11$--$0.17$ against $0.64$--$0.65$ when loaded, and a
no-load current ratio inside the motor band $[0.25,0.50]$. That is the
textbook signature of an energised induction motor drawing magnetising
current, recovered independently on four machines from unlabelled data,
and it is the strongest evidence in this paper that the state definition
is physical rather than fitted.

\textbf{On the fans and contactors there is no energised-idle state to
find, and the battery says so.} Their ``STANDBY'' components sit at
$0.2$--$9.3$\,W drawing $0.09$--$0.14$\,A: the model is splitting the
OFF band in two because it was told to find four states. The current
ratio is $0.002$--$0.009$ against an acceptance band starting at
$0.25$ on all four, and contactor~II additionally shows a
\emph{negative} power-factor separation. \textbf{This is the physics
battery doing the job it was built for.} The framework's state
definition is specific to motor-driven loads, and on the machines where
it does not apply it reports that rather than returning four confident
states regardless.

Two mechanisms behind those failures are worth naming, because both are
limitations of the method rather than of the data.

\emph{Ranking states by active power can misname them.} Exhaust fan~I's
``OFF'' component ($0.8\%$ of the record) carries $6.6$\,A at $0$\,W and
power factor $0$ --- standing reactive current with no real power ---
while its ``STANDBY'' component ($36.4\%$) sits at $0.3$\,W and
$0.13$\,A and is the machine's true de-energised state. The canonical
names are assigned by mean active power, so a near-zero-power component
carrying reactive current ranks below the genuine off state and takes
its name. No hour total changes, since both are non-productive, but it
is why the current-ratio check fails in the direction it does. The same
machine returns the largest forbidden-transition probability in the
plant ($3.4\times10^{-3}$, against $1.9\times10^{-4}$ on the four motor
machines), which is an independent measurement arriving at the same
conclusion (Section~\ref{sec:explain:transitions}).

\emph{The independent classifier fails where the model has nothing to
lock onto.} Otsu-thresholded power factor agrees with the state model on
$85$--$99.7\%$ of energised rows on seven machines and on $\mathbf{55.1\%}$
on exhaust fan~II --- the machine whose physics checks fail hardest.
Method~1's usefulness as an external check is thereby confirmed
\emph{and} its failure mode identified.

\subsection{Two results that generalise}
\label{sec:multimachine:general}

\textbf{The minimum-dwell result holds on all eight records.} Flicker
before the constraint ranges from $26.7\%$ (milling~II) to $77.7\%$
(exhaust fan~II), median $50.6\%$; after it, $0.00\%$ on seven machines
and $0.17\%$ on milling~I, at a cost of relabelling $0.15$--$6.1\%$ of
rows. What Section~\ref{sec:results:task4} established on one machine
--- that a transition prior cannot control flicker and an explicit
dwell constraint can --- now holds across duty cycles differing by an
order of magnitude.

\textbf{BIC does not select a state count here.} Widening the scan to
$k\in[2,8]$ returns $k=8$ on seven machines and $k=7$ on pelletizer~II,
with the criterion still falling at the boundary in seven of eight
cases; the improvement from $k=4$ to $k=8$ is $3.1$--$33.1\%$ of the
criterion's own value. This is a result about the criterion, not the
data (Section~\ref{sec:rw:states}). Fixing $k=4$ everywhere for
comparability is a choice, and this is the honest statement of what it
costs: \textbf{the physics vetoes are not a belt-and-braces addition on
top of BIC --- on records of this size they are what is actually doing
the selecting.}

\subsection{Which machines are hard, and why}
\label{sec:multimachine:difficulty}

Difficulty is measured on three axes: separability of the no-load/load
boundary (Cohen's $d$ and Bhattacharyya distance on log power), record
quality (coverage, spike rate, segment count), and agreement with the
independent classifier.

The ordering is not the one predicted. The expectation was that the
low-power auxiliary machines would be hardest. What the data show is
that the fans and contactors have the \emph{largest} separability
figures in the plant --- Cohen's $d$ of $9.3$ to $75.0$ against the
pelletizers' $4.3$ to $7.2$ --- while failing the physics checks,
because the separation being measured there is OFF against running, not
no-load against loaded. \textbf{Separability and validity are close to
orthogonal on this data}, which is the argument for keeping an external
physics battery rather than an internal cluster-quality score:
silhouette, BIC and Cohen's $d$ are all maximised by the partitions that
are physically wrong.

The two milling machines are hard for an unrelated reason. Their records
cover $12.2$ of $43$ days in $1\,082$ and $580$ fragments, against
$66$--$74$ fragments for the five-month records. They support two
one-week bootstrap blocks, which is not enough to form a usable
interval, and the intervals in Table~\ref{tab:multimachine} are printed
for those rows only to make that visible.
\textbf{Their $5.5$\,h/day --- the largest per-day standby figure among
the motor-driven machines --- must not be annualised}, and
Figure~\ref{fig:standby-hours} prints the block count inside each bar so
the figure cannot be read as the largest opportunity in the plant.

\begin{figure}[t]
  \centering
  \includegraphics[width=\linewidth]{fig_p6_standby_hours}
  \caption{STANDBY hours per covered day with $95\%$ weekly-block
  intervals. Bars are marked according to whether the machine passes the
  physics battery, and the number of weekly blocks is printed inside
  each bar; hatching marks an interval that rests on too few blocks to
  be usable.}
  \label{fig:standby-hours}
\end{figure}

\subsection{The decision layer exists on three machines of eight}
\label{sec:multimachine:decision}

A decision problem requires a measurable standby power and an observable
cold restart. Five machines fail an explicit degeneracy test rather than
an inspection: four because standby power is below $1$\,W, so the
standby cost per second is numerically zero, the break-even point is
infinite, and every currency figure would be an artefact of dividing by
zero; and pelletizer~II because in five months it is \emph{never once
observed resuming production from OFF}, so the cold-versus-warm contrast
cannot be measured at all. Two machines return a \emph{negative} restart
delay --- cold starts reaching production faster than warm ones --- which
is not a plausible physical reading and is a further symptom of the same
problem. The three machines that survive the test are the subject of
Table~\ref{tab:multimachine:decision}.

\begin{table}[t]
\centering
\caption{Machines on which the shutdown decision is well posed.
Head-room is the oracle saving; the proposed policy's figure is quoted
against the plant's own behaviour.}
\label{tab:multimachine:decision}
\small
\setlength{\tabcolsep}{4pt}
\begin{tabular}{@{}lrrrr@{}}
\toprule
Machine & energised idle & head-room & forecast MAE & proposed \\
        & (h, \% of idle) & (USD/yr) & (s) & (USD/yr) \\
\midrule
Pelletizer I & $88$ ($18.5\%$) & $37.6$ & $9\,034$ & $+2.9$ \\
Milling I    & $67$ ($48.8\%$) & $24.9$ & $18\,205$ & $\mathbf{-1\,041}$ \\
Milling II   & $64$ ($47.6\%$) & $25.2$ & $18\,440$ & $\mathbf{-1\,256}$ \\
\bottomrule
\end{tabular}
\end{table}

\textbf{The conjecture that motivated this experiment is refuted.} The
open question after the single-machine study was whether some other
IMDELD machine has enough recoverable standby to fund a retrofit ---
higher standby power, a larger energised-idle fraction, or a plant that
does not already shut down manually. None does. The three machines with
a well-posed problem have head-rooms of $37.6$, $24.9$ and
$25.2$\,USD/yr, so \textbf{the facility's total attainable saving from
standby recovery, summed over every machine where the question is even
meaningful, is under $90$\,USD/yr}. The framework's contribution on this
facility is decision quality, not recovered energy, and this paper says
so in its abstract rather than at the end of its discussion.

\textbf{The milling machines are worse than that, and the negative
result is reported as it stands.} The proposed policy \emph{loses}
$1\,041$ and $1\,256$\,USD/yr there against a head-room of
$25$\,USD/yr. The cause is visible in the same table: a duration
forecaster with MAE $18\,200$\,s and Spearman $0.34$--$0.42$, against
$9\,034$\,s and $0.58$ on pelletizer~I, fitted on a record broken into
$1\,082$ fragments. It over-predicts remaining idle time and fires
shutdowns the realised episode does not support. The decision layer
inherits its forecaster's quality, and on a record this fragmented the
forecaster is not good enough to act on --- which is a deployment
criterion, not merely an observation.

%  Figures are referenced by bare name; add to the preamble
%      \graphicspath{{../outputs/phase4/figures/}{../outputs/phase6/figures/}}
%  Bibliography: paper/references.bib
%  Numbers below are generated by  python run_phase4.py --only task8
%  and stored in outputs/phase4/task8/table_*.csv, with the intervals
%  from outputs/phase5/task11/state_block_bootstrap.csv.
% =====================================================================

\section{The Eight Machines}
\label{sec:multimachine}

Everything to this point is pelletizer~I. This section applies the
identical procedure --- one preprocessing, one labeller, $k=4$ fixed,
one physics battery --- to all eight IMDELD machines, and asks two
questions of the result: on which machines is the state definition
physically valid, and on which is there a shutdown decision to make at
all. The two answers are not the same, and neither follows the
prediction that motivated the experiment.

\subsection{Uniform labelling}
\label{sec:multimachine:labelling}

The Phase~I--III labels exist only for pelletizer~I and were produced by
a script whose configuration has since moved, so comparing that export
against seven fresh labellings would confound machine differences with
procedure differences. All eight records are therefore re-labelled
identically. Three departures from the original implementation were
required, each verified rather than assumed.

\emph{Blank currents.} The two milling-machine meters leave the current
channel empty on $36.1\%$ and $36.3\%$ of rows --- matching their OFF
shares almost exactly, i.e.\ the meter blanks the channel whenever the
machine is de-energised. On every such row the apparent power is exactly
$0$\,VA with voltage present, and $S=VI$ then admits only $I=0$\,A, so
the cells are filled with zero and a unit test asserts that every filled
cell met that condition. Leaving them missing makes a GMM
unfittable; dropping them would delete a third of each record precisely
where the machine is off, biasing every state-time total.

\emph{Timestamps and storage.} IMDELD spans a daylight-saving change, so
the timestamp column carries mixed UTC offsets; splitting the naive part
from the offset is exact and ten times faster than the fallback parse.
Records are stored as \texttt{float32}, whose one measurable consequence
is that about two rows per $10^{5}$ fall on the other side of the
spike threshold.

On pelletizer~I --- the one machine both labellings cover --- the
uniform labeller reproduces the pipeline's own export at $98.88\%$ row
agreement, $99.84\%$ on STANDBY-versus-rest, ARI $0.978$, and $89$\,h of
STANDBY against the export's $87$\,h. A number in this section that
differs from an earlier one is therefore a consequence of the analysis,
not of the labels.

Both of those figures are single-seed. The total this paper quotes for
pelletizer~I is the physics-conditional mean over passing seeds,
$92\pm3$\,h (Section~\ref{sec:method:accept_reject}); the single-seed
values sit inside that interval, and the per-machine rows of
Table~\ref{tab:multimachine} are reported per covered day with a
sampling interval for the same reason.

\subsection{Where the state definition holds}
\label{sec:multimachine:states}

\begin{table*}[!t]
\centering
\caption{Per-machine characterisation. STANDBY hours per covered day
carry a $95\%$ one-week moving-block bootstrap interval
(Section~\ref{sec:significance}); $n_{b}$ is the number of blocks the
record supports. $P_{\text{sb}}$ is the mean power of the inferred
STANDBY component. The no-load current ratio has an acceptance band of
$[0.25,0.50]$ from induction-motor theory. M1 is the fraction of
energised rows on which the independent Otsu power-factor classifier
agrees with the state model.}
\label{tab:multimachine}
\small
\begin{tabular}{@{}lrrrrrrrr@{}}
\toprule
Machine & cov.\ (d) & STANDBY (h/d) & $n_{b}$ & $P_{\text{sb}}$ & physics & PF sep. & $I$ ratio & M1 \\
\midrule
\multicolumn{9}{@{}l}{\textit{Motor-driven --- the state definition applies}}\\
Pelletizer I    & $63.4$ & $1.40$\,{\footnotesize$[1.15,1.63]$} & $10$ & $4.2$\,kW & $\mathbf{5/5}$ & $0.72$ & $0.36$ & $99.4\%$ \\
Pelletizer II   & $62.4$ & $1.85$\,{\footnotesize$[1.47,2.28]$} & $9$  & $4.5$\,kW & $\mathbf{5/5}$ & $0.71$ & $0.36$ & $99.7\%$ \\
Milling I       & $12.2$ & $5.52$\,{\footnotesize$[4.58,6.23]$} & $\mathbf{2}$ & $4.6$\,kW & $\mathbf{5/5}$ & $0.61$ & $0.43$ & $98.1\%$ \\
Milling II      & $12.2$ & $5.27$\,{\footnotesize$[4.44,5.89]$} & $\mathbf{2}$ & $5.4$\,kW & $\mathbf{5/5}$ & $0.52$ & $0.33$ & $97.7\%$ \\
\addlinespace
\multicolumn{9}{@{}l}{\textit{Auxiliary --- no energised-idle state exists to find}}\\
Exhaust fan I   & $63.7$ & $8.73$\,{\footnotesize$[7.31,10.03]$} & $10$ & $0.3$\,W & $4/5$ & $0.91$ & $\mathbf{0.008}$ & $96.6\%$ \\
Exhaust fan II  & $62.3$ & $1.45$\,{\footnotesize$[1.09,2.07]$}  & $9$  & $0.2$\,W & $3/5$ & $0.62$ & $\mathbf{0.002}$ & $\mathbf{55.1\%}$ \\
Contactor I     & $63.7$ & $0.90$\,{\footnotesize$[0.71,1.12]$}  & $10$ & $0.3$\,W & $3/5$ & $0.31$ & $\mathbf{0.009}$ & $90.3\%$ \\
Contactor II    & $62.5$ & $1.15$\,{\footnotesize$[0.85,1.42]$}  & $9$  & $9.3$\,W & $3/5$ & $\mathbf{-0.05}$ & $\mathbf{0.007}$ & $85.4\%$ \\
\bottomrule
\end{tabular}
\end{table*}

\textbf{The two pelletizers and the two milling machines pass all five
physics checks.} Their STANDBY components sit at $4.2$--$5.4$\,kW with a
power factor of $0.11$--$0.17$ against $0.64$--$0.65$ when loaded, and a
no-load current ratio inside the motor band $[0.25,0.50]$. That is the
textbook signature of an energised induction motor drawing magnetising
current, recovered independently on four machines from unlabelled data,
and it is the strongest evidence in this paper that the state definition
is physical rather than fitted.

\textbf{On the fans and contactors there is no energised-idle state to
find, and the battery says so.} Their ``STANDBY'' components sit at
$0.2$--$9.3$\,W drawing $0.09$--$0.14$\,A: the model is splitting the
OFF band in two because it was told to find four states. The current
ratio is $0.002$--$0.009$ against an acceptance band starting at
$0.25$ on all four, and contactor~II additionally shows a
\emph{negative} power-factor separation. \textbf{This is the physics
battery doing the job it was built for.} The framework's state
definition is specific to motor-driven loads, and on the machines where
it does not apply it reports that rather than returning four confident
states regardless.

Two mechanisms behind those failures are worth naming, because both are
limitations of the method rather than of the data.

\emph{Ranking states by active power can misname them.} Exhaust fan~I's
``OFF'' component ($0.8\%$ of the record) carries $6.6$\,A at $0$\,W and
power factor $0$ --- standing reactive current with no real power ---
while its ``STANDBY'' component ($36.4\%$) sits at $0.3$\,W and
$0.13$\,A and is the machine's true de-energised state. The canonical
names are assigned by mean active power, so a near-zero-power component
carrying reactive current ranks below the genuine off state and takes
its name. No hour total changes, since both are non-productive, but it
is why the current-ratio check fails in the direction it does. The same
machine returns the largest forbidden-transition probability in the
plant ($3.4\times10^{-3}$, against $1.9\times10^{-4}$ on the four motor
machines), which is an independent measurement arriving at the same
conclusion (Section~\ref{sec:explain:transitions}).

\emph{The independent classifier fails where the model has nothing to
lock onto.} Otsu-thresholded power factor agrees with the state model on
$85$--$99.7\%$ of energised rows on seven machines and on $\mathbf{55.1\%}$
on exhaust fan~II --- the machine whose physics checks fail hardest.
Method~1's usefulness as an external check is thereby confirmed
\emph{and} its failure mode identified.

\subsection{Two results that generalise}
\label{sec:multimachine:general}

\textbf{The minimum-dwell result holds on all eight records.} Flicker
before the constraint ranges from $26.7\%$ (milling~II) to $77.7\%$
(exhaust fan~II), median $50.6\%$; after it, $0.00\%$ on seven machines
and $0.17\%$ on milling~I, at a cost of relabelling $0.15$--$6.1\%$ of
rows. What Section~\ref{sec:results:task4} established on one machine
--- that a transition prior cannot control flicker and an explicit
dwell constraint can --- now holds across duty cycles differing by an
order of magnitude.

\textbf{BIC does not select a state count here.} Widening the scan to
$k\in[2,8]$ returns $k=8$ on seven machines and $k=7$ on pelletizer~II,
with the criterion still falling at the boundary in seven of eight
cases; the improvement from $k=4$ to $k=8$ is $3.1$--$33.1\%$ of the
criterion's own value. This is a result about the criterion, not the
data (Section~\ref{sec:rw:states}). Fixing $k=4$ everywhere for
comparability is a choice, and this is the honest statement of what it
costs: \textbf{the physics vetoes are not a belt-and-braces addition on
top of BIC --- on records of this size they are what is actually doing
the selecting.}

\subsection{Which machines are hard, and why}
\label{sec:multimachine:difficulty}

Difficulty is measured on three axes: separability of the no-load/load
boundary (Cohen's $d$ and Bhattacharyya distance on log power), record
quality (coverage, spike rate, segment count), and agreement with the
independent classifier.

The ordering is not the one predicted. The expectation was that the
low-power auxiliary machines would be hardest. What the data show is
that the fans and contactors have the \emph{largest} separability
figures in the plant --- Cohen's $d$ of $9.3$ to $75.0$ against the
pelletizers' $4.3$ to $7.2$ --- while failing the physics checks,
because the separation being measured there is OFF against running, not
no-load against loaded. \textbf{Separability and validity are close to
orthogonal on this data}, which is the argument for keeping an external
physics battery rather than an internal cluster-quality score:
silhouette, BIC and Cohen's $d$ are all maximised by the partitions that
are physically wrong.

The two milling machines are hard for an unrelated reason. Their records
cover $12.2$ of $43$ days in $1\,082$ and $580$ fragments, against
$66$--$74$ fragments for the five-month records. They support two
one-week bootstrap blocks, which is not enough to form a usable
interval, and the intervals in Table~\ref{tab:multimachine} are printed
for those rows only to make that visible.
\textbf{Their $5.5$\,h/day --- the largest per-day standby figure among
the motor-driven machines --- must not be annualised}, and
Figure~\ref{fig:standby-hours} prints the block count inside each bar so
the figure cannot be read as the largest opportunity in the plant.

\begin{figure}[t]
  \centering
  \includegraphics[width=\linewidth]{fig_p6_standby_hours}
  \caption{STANDBY hours per covered day with $95\%$ weekly-block
  intervals. Bars are marked according to whether the machine passes the
  physics battery, and the number of weekly blocks is printed inside
  each bar; hatching marks an interval that rests on too few blocks to
  be usable.}
  \label{fig:standby-hours}
\end{figure}

\subsection{The decision layer exists on three machines of eight}
\label{sec:multimachine:decision}

A decision problem requires a measurable standby power and an observable
cold restart. Five machines fail an explicit degeneracy test rather than
an inspection: four because standby power is below $1$\,W, so the
standby cost per second is numerically zero, the break-even point is
infinite, and every currency figure would be an artefact of dividing by
zero; and pelletizer~II because in five months it is \emph{never once
observed resuming production from OFF}, so the cold-versus-warm contrast
cannot be measured at all. Two machines return a \emph{negative} restart
delay --- cold starts reaching production faster than warm ones --- which
is not a plausible physical reading and is a further symptom of the same
problem. The three machines that survive the test are the subject of
Table~\ref{tab:multimachine:decision}.

\begin{table}[t]
\centering
\caption{Machines on which the shutdown decision is well posed.
Head-room is the oracle saving; the proposed policy's figure is quoted
against the plant's own behaviour.}
\label{tab:multimachine:decision}
\small
\setlength{\tabcolsep}{4pt}
\begin{tabular}{@{}lrrrr@{}}
\toprule
Machine & energised idle & head-room & forecast MAE & proposed \\
        & (h, \% of idle) & (USD/yr) & (s) & (USD/yr) \\
\midrule
Pelletizer I & $88$ ($18.5\%$) & $37.6$ & $9\,034$ & $+2.9$ \\
Milling I    & $67$ ($48.8\%$) & $24.9$ & $18\,205$ & $\mathbf{-1\,041}$ \\
Milling II   & $64$ ($47.6\%$) & $25.2$ & $18\,440$ & $\mathbf{-1\,256}$ \\
\bottomrule
\end{tabular}
\end{table}

\textbf{The conjecture that motivated this experiment is refuted.} The
open question after the single-machine study was whether some other
IMDELD machine has enough recoverable standby to fund a retrofit ---
higher standby power, a larger energised-idle fraction, or a plant that
does not already shut down manually. None does. The three machines with
a well-posed problem have head-rooms of $37.6$, $24.9$ and
$25.2$\,USD/yr, so \textbf{the facility's total attainable saving from
standby recovery, summed over every machine where the question is even
meaningful, is under $90$\,USD/yr}. The framework's contribution on this
facility is decision quality, not recovered energy, and this paper says
so in its abstract rather than at the end of its discussion.

\textbf{The milling machines are worse than that, and the negative
result is reported as it stands.} The proposed policy \emph{loses}
$1\,041$ and $1\,256$\,USD/yr there against a head-room of
$25$\,USD/yr. The cause is visible in the same table: a duration
forecaster with MAE $18\,200$\,s and Spearman $0.34$--$0.42$, against
$9\,034$\,s and $0.58$ on pelletizer~I, fitted on a record broken into
$1\,082$ fragments. It over-predicts remaining idle time and fires
shutdowns the realised episode does not support. The decision layer
inherits its forecaster's quality, and on a record this fragmented the
forecaster is not good enough to act on --- which is a deployment
criterion, not merely an observation.

%  Figures are referenced by bare name; add to the preamble
%      \graphicspath{{../outputs/phase4/figures/}{../outputs/phase6/figures/}}
%  Bibliography: paper/references.bib
%  Numbers below are generated by  python run_phase4.py --only task8
%  and stored in outputs/phase4/task8/table_*.csv, with the intervals
%  from outputs/phase5/task11/state_block_bootstrap.csv.
% =====================================================================

\section{The Eight Machines}
\label{sec:multimachine}

Everything to this point is pelletizer~I. This section applies the
identical procedure --- one preprocessing, one labeller, $k=4$ fixed,
one physics battery --- to all eight IMDELD machines, and asks two
questions of the result: on which machines is the state definition
physically valid, and on which is there a shutdown decision to make at
all. The two answers are not the same, and neither follows the
prediction that motivated the experiment.

\subsection{Uniform labelling}
\label{sec:multimachine:labelling}

The Phase~I--III labels exist only for pelletizer~I and were produced by
a script whose configuration has since moved, so comparing that export
against seven fresh labellings would confound machine differences with
procedure differences. All eight records are therefore re-labelled
identically. Three departures from the original implementation were
required, each verified rather than assumed.

\emph{Blank currents.} The two milling-machine meters leave the current
channel empty on $36.1\%$ and $36.3\%$ of rows --- matching their OFF
shares almost exactly, i.e.\ the meter blanks the channel whenever the
machine is de-energised. On every such row the apparent power is exactly
$0$\,VA with voltage present, and $S=VI$ then admits only $I=0$\,A, so
the cells are filled with zero and a unit test asserts that every filled
cell met that condition. Leaving them missing makes a GMM
unfittable; dropping them would delete a third of each record precisely
where the machine is off, biasing every state-time total.

\emph{Timestamps and storage.} IMDELD spans a daylight-saving change, so
the timestamp column carries mixed UTC offsets; splitting the naive part
from the offset is exact and ten times faster than the fallback parse.
Records are stored as \texttt{float32}, whose one measurable consequence
is that about two rows per $10^{5}$ fall on the other side of the
spike threshold.

On pelletizer~I --- the one machine both labellings cover --- the
uniform labeller reproduces the pipeline's own export at $98.88\%$ row
agreement, $99.84\%$ on STANDBY-versus-rest, ARI $0.978$, and $89$\,h of
STANDBY against the export's $87$\,h. A number in this section that
differs from an earlier one is therefore a consequence of the analysis,
not of the labels.

Both of those figures are single-seed. The total this paper quotes for
pelletizer~I is the physics-conditional mean over passing seeds,
$92\pm3$\,h (Section~\ref{sec:method:accept_reject}); the single-seed
values sit inside that interval, and the per-machine rows of
Table~\ref{tab:multimachine} are reported per covered day with a
sampling interval for the same reason.

\subsection{Where the state definition holds}
\label{sec:multimachine:states}

\begin{table*}[!t]
\centering
\caption{Per-machine characterisation. STANDBY hours per covered day
carry a $95\%$ one-week moving-block bootstrap interval
(Section~\ref{sec:significance}); $n_{b}$ is the number of blocks the
record supports. $P_{\text{sb}}$ is the mean power of the inferred
STANDBY component. The no-load current ratio has an acceptance band of
$[0.25,0.50]$ from induction-motor theory. M1 is the fraction of
energised rows on which the independent Otsu power-factor classifier
agrees with the state model.}
\label{tab:multimachine}
\small
\begin{tabular}{@{}lrrrrrrrr@{}}
\toprule
Machine & cov.\ (d) & STANDBY (h/d) & $n_{b}$ & $P_{\text{sb}}$ & physics & PF sep. & $I$ ratio & M1 \\
\midrule
\multicolumn{9}{@{}l}{\textit{Motor-driven --- the state definition applies}}\\
Pelletizer I    & $63.4$ & $1.40$\,{\footnotesize$[1.15,1.63]$} & $10$ & $4.2$\,kW & $\mathbf{5/5}$ & $0.72$ & $0.36$ & $99.4\%$ \\
Pelletizer II   & $62.4$ & $1.85$\,{\footnotesize$[1.47,2.28]$} & $9$  & $4.5$\,kW & $\mathbf{5/5}$ & $0.71$ & $0.36$ & $99.7\%$ \\
Milling I       & $12.2$ & $5.52$\,{\footnotesize$[4.58,6.23]$} & $\mathbf{2}$ & $4.6$\,kW & $\mathbf{5/5}$ & $0.61$ & $0.43$ & $98.1\%$ \\
Milling II      & $12.2$ & $5.27$\,{\footnotesize$[4.44,5.89]$} & $\mathbf{2}$ & $5.4$\,kW & $\mathbf{5/5}$ & $0.52$ & $0.33$ & $97.7\%$ \\
\addlinespace
\multicolumn{9}{@{}l}{\textit{Auxiliary --- no energised-idle state exists to find}}\\
Exhaust fan I   & $63.7$ & $8.73$\,{\footnotesize$[7.31,10.03]$} & $10$ & $0.3$\,W & $4/5$ & $0.91$ & $\mathbf{0.008}$ & $96.6\%$ \\
Exhaust fan II  & $62.3$ & $1.45$\,{\footnotesize$[1.09,2.07]$}  & $9$  & $0.2$\,W & $3/5$ & $0.62$ & $\mathbf{0.002}$ & $\mathbf{55.1\%}$ \\
Contactor I     & $63.7$ & $0.90$\,{\footnotesize$[0.71,1.12]$}  & $10$ & $0.3$\,W & $3/5$ & $0.31$ & $\mathbf{0.009}$ & $90.3\%$ \\
Contactor II    & $62.5$ & $1.15$\,{\footnotesize$[0.85,1.42]$}  & $9$  & $9.3$\,W & $3/5$ & $\mathbf{-0.05}$ & $\mathbf{0.007}$ & $85.4\%$ \\
\bottomrule
\end{tabular}
\end{table*}

\textbf{The two pelletizers and the two milling machines pass all five
physics checks.} Their STANDBY components sit at $4.2$--$5.4$\,kW with a
power factor of $0.11$--$0.17$ against $0.64$--$0.65$ when loaded, and a
no-load current ratio inside the motor band $[0.25,0.50]$. That is the
textbook signature of an energised induction motor drawing magnetising
current, recovered independently on four machines from unlabelled data,
and it is the strongest evidence in this paper that the state definition
is physical rather than fitted.

\textbf{On the fans and contactors there is no energised-idle state to
find, and the battery says so.} Their ``STANDBY'' components sit at
$0.2$--$9.3$\,W drawing $0.09$--$0.14$\,A: the model is splitting the
OFF band in two because it was told to find four states. The current
ratio is $0.002$--$0.009$ against an acceptance band starting at
$0.25$ on all four, and contactor~II additionally shows a
\emph{negative} power-factor separation. \textbf{This is the physics
battery doing the job it was built for.} The framework's state
definition is specific to motor-driven loads, and on the machines where
it does not apply it reports that rather than returning four confident
states regardless.

Two mechanisms behind those failures are worth naming, because both are
limitations of the method rather than of the data.

\emph{Ranking states by active power can misname them.} Exhaust fan~I's
``OFF'' component ($0.8\%$ of the record) carries $6.6$\,A at $0$\,W and
power factor $0$ --- standing reactive current with no real power ---
while its ``STANDBY'' component ($36.4\%$) sits at $0.3$\,W and
$0.13$\,A and is the machine's true de-energised state. The canonical
names are assigned by mean active power, so a near-zero-power component
carrying reactive current ranks below the genuine off state and takes
its name. No hour total changes, since both are non-productive, but it
is why the current-ratio check fails in the direction it does. The same
machine returns the largest forbidden-transition probability in the
plant ($3.4\times10^{-3}$, against $1.9\times10^{-4}$ on the four motor
machines), which is an independent measurement arriving at the same
conclusion (Section~\ref{sec:explain:transitions}).

\emph{The independent classifier fails where the model has nothing to
lock onto.} Otsu-thresholded power factor agrees with the state model on
$85$--$99.7\%$ of energised rows on seven machines and on $\mathbf{55.1\%}$
on exhaust fan~II --- the machine whose physics checks fail hardest.
Method~1's usefulness as an external check is thereby confirmed
\emph{and} its failure mode identified.

\subsection{Two results that generalise}
\label{sec:multimachine:general}

\textbf{The minimum-dwell result holds on all eight records.} Flicker
before the constraint ranges from $26.7\%$ (milling~II) to $77.7\%$
(exhaust fan~II), median $50.6\%$; after it, $0.00\%$ on seven machines
and $0.17\%$ on milling~I, at a cost of relabelling $0.15$--$6.1\%$ of
rows. What Section~\ref{sec:results:task4} established on one machine
--- that a transition prior cannot control flicker and an explicit
dwell constraint can --- now holds across duty cycles differing by an
order of magnitude.

\textbf{BIC does not select a state count here.} Widening the scan to
$k\in[2,8]$ returns $k=8$ on seven machines and $k=7$ on pelletizer~II,
with the criterion still falling at the boundary in seven of eight
cases; the improvement from $k=4$ to $k=8$ is $3.1$--$33.1\%$ of the
criterion's own value. This is a result about the criterion, not the
data (Section~\ref{sec:rw:states}). Fixing $k=4$ everywhere for
comparability is a choice, and this is the honest statement of what it
costs: \textbf{the physics vetoes are not a belt-and-braces addition on
top of BIC --- on records of this size they are what is actually doing
the selecting.}

\subsection{Which machines are hard, and why}
\label{sec:multimachine:difficulty}

Difficulty is measured on three axes: separability of the no-load/load
boundary (Cohen's $d$ and Bhattacharyya distance on log power), record
quality (coverage, spike rate, segment count), and agreement with the
independent classifier.

The ordering is not the one predicted. The expectation was that the
low-power auxiliary machines would be hardest. What the data show is
that the fans and contactors have the \emph{largest} separability
figures in the plant --- Cohen's $d$ of $9.3$ to $75.0$ against the
pelletizers' $4.3$ to $7.2$ --- while failing the physics checks,
because the separation being measured there is OFF against running, not
no-load against loaded. \textbf{Separability and validity are close to
orthogonal on this data}, which is the argument for keeping an external
physics battery rather than an internal cluster-quality score:
silhouette, BIC and Cohen's $d$ are all maximised by the partitions that
are physically wrong.

The two milling machines are hard for an unrelated reason. Their records
cover $12.2$ of $43$ days in $1\,082$ and $580$ fragments, against
$66$--$74$ fragments for the five-month records. They support two
one-week bootstrap blocks, which is not enough to form a usable
interval, and the intervals in Table~\ref{tab:multimachine} are printed
for those rows only to make that visible.
\textbf{Their $5.5$\,h/day --- the largest per-day standby figure among
the motor-driven machines --- must not be annualised}, and
Figure~\ref{fig:standby-hours} prints the block count inside each bar so
the figure cannot be read as the largest opportunity in the plant.

\begin{figure}[t]
  \centering
  \includegraphics[width=\linewidth]{fig_p6_standby_hours}
  \caption{STANDBY hours per covered day with $95\%$ weekly-block
  intervals. Bars are marked according to whether the machine passes the
  physics battery, and the number of weekly blocks is printed inside
  each bar; hatching marks an interval that rests on too few blocks to
  be usable.}
  \label{fig:standby-hours}
\end{figure}

\subsection{The decision layer exists on three machines of eight}
\label{sec:multimachine:decision}

A decision problem requires a measurable standby power and an observable
cold restart. Five machines fail an explicit degeneracy test rather than
an inspection: four because standby power is below $1$\,W, so the
standby cost per second is numerically zero, the break-even point is
infinite, and every currency figure would be an artefact of dividing by
zero; and pelletizer~II because in five months it is \emph{never once
observed resuming production from OFF}, so the cold-versus-warm contrast
cannot be measured at all. Two machines return a \emph{negative} restart
delay --- cold starts reaching production faster than warm ones --- which
is not a plausible physical reading and is a further symptom of the same
problem. The three machines that survive the test are the subject of
Table~\ref{tab:multimachine:decision}.

\begin{table}[t]
\centering
\caption{Machines on which the shutdown decision is well posed.
Head-room is the oracle saving; the proposed policy's figure is quoted
against the plant's own behaviour.}
\label{tab:multimachine:decision}
\small
\setlength{\tabcolsep}{4pt}
\begin{tabular}{@{}lrrrr@{}}
\toprule
Machine & energised idle & head-room & forecast MAE & proposed \\
        & (h, \% of idle) & (USD/yr) & (s) & (USD/yr) \\
\midrule
Pelletizer I & $88$ ($18.5\%$) & $37.6$ & $9\,034$ & $+2.9$ \\
Milling I    & $67$ ($48.8\%$) & $24.9$ & $18\,205$ & $\mathbf{-1\,041}$ \\
Milling II   & $64$ ($47.6\%$) & $25.2$ & $18\,440$ & $\mathbf{-1\,256}$ \\
\bottomrule
\end{tabular}
\end{table}

\textbf{The conjecture that motivated this experiment is refuted.} The
open question after the single-machine study was whether some other
IMDELD machine has enough recoverable standby to fund a retrofit ---
higher standby power, a larger energised-idle fraction, or a plant that
does not already shut down manually. None does. The three machines with
a well-posed problem have head-rooms of $37.6$, $24.9$ and
$25.2$\,USD/yr, so \textbf{the facility's total attainable saving from
standby recovery, summed over every machine where the question is even
meaningful, is under $90$\,USD/yr}. The framework's contribution on this
facility is decision quality, not recovered energy, and this paper says
so in its abstract rather than at the end of its discussion.

\textbf{The milling machines are worse than that, and the negative
result is reported as it stands.} The proposed policy \emph{loses}
$1\,041$ and $1\,256$\,USD/yr there against a head-room of
$25$\,USD/yr. The cause is visible in the same table: a duration
forecaster with MAE $18\,200$\,s and Spearman $0.34$--$0.42$, against
$9\,034$\,s and $0.58$ on pelletizer~I, fitted on a record broken into
$1\,082$ fragments. It over-predicts remaining idle time and fires
shutdowns the realised episode does not support. The decision layer
inherits its forecaster's quality, and on a record this fragmented the
forecaster is not good enough to act on --- which is a deployment
criterion, not merely an observation.
